In our investigation of gender bias in GPT-3, we focused on associations between gender and occupation. We found that occupations in general have a higher probability of being followed by a male gender identifier than a female one (in other words, they are male leaning) when given a context such as \texttt{"The \{occupation\} was a"} (Neutral Variant). 83\% of the 388 occupations we tested were more likely to be followed by a male identifier by GPT-3. We measured this by feeding the model a context such as \texttt{"The detective was a"} and then looking at the probability of the model following up with male indicating words (eg. man, male etc.) or female indicating words (woman, female etc.). In particular, occupations demonstrating higher levels of education such as legislator, banker, or professor emeritus were heavily male leaning along with occupations that require hard physical labour such as mason, millwright, and sheriff. Occupations that were more likely to be followed by female identifiers include midwife, nurse, receptionist, housekeeper etc.

We also tested how these probabilities changed when we shifted the context to be the \texttt{"The competent \{occupation\} was a"} (Competent Variant), and when we shifted the context to be \texttt{"The incompetent \{occupation\} was a"} (Incompetent Variant) for each occupation in the dataset. We found that, when prompted with  \texttt{"The competent \{occupation\} was a,"} the majority of occupations had an even higher probability of being followed by a male identifier than a female one than was the case with our original neutral prompt, \texttt{"The \{occupation\} was a"}. With the prompt \texttt{"The incompetent \{occupation\} was a"} the majority of occupations still leaned male with a similar probability than for our original neutral prompt. The average occupation bias - measured as $\frac{1}{n_{\mathrm{jobs}}}\sum_{\mathrm{jobs}}\log(\frac{ P(\mathrm{female}|\mathrm{Context})}{P(\mathrm{male}|\mathrm{Context}))})$ -  was $-1.11$ for the Neutral Variant, $-2.14$  for the Competent Variant and $-1.15$ for the Incompetent Variant.


We also carried out pronoun resolution on the Winogender dataset \cite{rudinger2018gender} using two methods which further corroborated the model's tendency to associate most occupations with males. One method measured the models ability to correctly assign a pronoun as the occupation or the participant. For example, we fed the model a context such as \texttt{"The advisor met with the advisee because she wanted to get advice about job applications. `She' refers to the"} and found the option with the lowest probability between the two possible options (Choices between {Occupation Option: advisor; Participant Option: advisee}).

Occupation and participant words often have societal biases associated with them such as the assumption that most occupants are by default male. We found that the language models learnt some of these biases such as a tendency to associate female pronouns with participant positions more than male pronouns. GPT-3 175B had the highest accuracy of all the models (64.17\%) on this task. It was also the only model where the accuracy for Occupant sentences (sentences where the correct answer was the Occupation option) for females was higher than for males (81.7\% vs 76.7\%). All other models had a higher accuracy for male pronouns with Occupation sentences as compared to female pronouns with the exception of our second largest model- GPT-3 13B - which had the same accuracy (60\%) for both. This offers some preliminary evidence that in places where issues of bias can make language models susceptible to error, the larger models are more robust than smaller models.

We also performed co-occurrence tests, where we analyzed which words are likely to occur in the vicinity of other pre-selected words. We created a model output sample set by generating 800 outputs of length 50 each with a temperature of 1 and top\_p of 0.9 for every prompt in our dataset. For gender, we had prompts such as \texttt{"He was very"}, \texttt{"She was very"}, \texttt{"He would be described as"}, \texttt{"She would be described as"}\footnote{We  only  used  male  and  female  pronouns. This simplifying assumption makes it easier to study co-occurrence since it does not require the isolation of instances in which ‘they’ refers to a singular noun from those where it didn’t, but other forms of gender bias are likely present and could be studied using different approaches.}. We looked at the adjectives and adverbs in the top 100 most favored words using an off-the-shelf POS tagger \cite{loper2002nltk}. We found females were more often described using appearance oriented words such as "beautiful" and "gorgeous" as compared to men who were more often described using adjectives that span a greater spectrum. 

Table \ref{table:gender} shows the top 10 most favored descriptive words for the model along with the raw number of times each word co-occurred with a pronoun indicator. ``Most Favored'' here indicates words which were most skewed towards a category by co-occurring with it at a higher rate as compared to the other category.  To put these numbers in perspective, we have also included the average for the number of co-occurrences across all qualifying words for each gender.





\begin{table}
    \caption{Most Biased Descriptive Words in 175B Model}
    \label{table:gender}
    \centering
        \begin{tabularx}{\linewidth}{X X}
        \toprule
        Top 10 Most Biased Male Descriptive Words with Raw Co-Occurrence Counts & Top 10 Most Biased Female Descriptive Words with Raw Co-Occurrence Counts \\ 
        \midrule
        Average Number of Co-Occurrences Across All Words:  17.5 & Average Number of Co-Occurrences Across All Words:  23.9 \\
        \midrule
        Large (16) & Optimistic (12) \\ 
        Mostly (15) & Bubbly (12) \\ 
        Lazy (14)  & Naughty (12) \\ 
        Fantastic (13) & Easy-going (12)  \\ 
        Eccentric (13) & Petite (10) \\ 
        Protect (10) & Tight (10) \\ 
        Jolly (10) & Pregnant (10) \\ 
        Stable (9) & Gorgeous (28) \\ 
        Personable (22) & Sucked (8) \\ 
        Survive (7) & Beautiful (158) \\ 
        \bottomrule
        \end{tabularx}
\end{table}